\documentclass[aspectratio=169, 10pt]{beamer}
\usetheme{Madrid}
\usecolortheme{default}

\usepackage[utf8]{inputenc}
\usepackage[T1]{fontenc}
\usepackage{lmodern}
\usepackage[french]{babel}
\usepackage{amsmath,amssymb,amsfonts}
\usepackage{booktabs}
\usepackage{colortbl}
\usepackage{xcolor}
\usepackage{tikz}
\usetikzlibrary{arrows.meta,positioning,shapes,calc}
\usepackage{pgfplots}
\pgfplotsset{compat=1.18}

\setbeamercovered{transparent}
\setbeamertemplate{navigation symbols}{}


\definecolor{geblue}{RGB}{31, 119, 180}
\definecolor{gebluemid}{RGB}{31, 119, 180}
\definecolor{cghred}{RGB}{214, 39, 40}
\definecolor{lightgeblue}{RGB}{165, 200, 224}
\definecolor{barlightblue}{RGB}{121, 173, 210}
\definecolor{dipgorange}{RGB}{255, 127, 14}
\definecolor{midlgreen}{RGB}{44, 160, 44}
\definecolor{midlbrown}{RGB}{147, 39, 39}
\definecolor{boxblue}{RGB}{221, 234, 243}
\definecolor{boxred}{RGB}{248, 222, 222}
\definecolor{boxgreen}{RGB}{223, 240, 223}
\definecolor{boxlavender}{RGB}{235, 237, 244}
\definecolor{cellgreen}{RGB}{171, 217, 171}
\definecolor{cellred}{RGB}{245, 201, 201}
\definecolor{celllightred}{RGB}{251, 233, 233}

\title[Localisation de tumeurs cérébrales]{Classification de la localisation \\ de tumeurs cérébrales pédiatriques}
\subtitle{Intégration multi-omique (expression génique + CGH) en régime $p \gg n$}
\author[]{R. Barata \quad A. De Rouck \quad M. Hoizey-Cangioloni \quad P. Lubaszewski Lima \quad M. Seyedi}
\institute[CentraleSupélec]{CentraleSupélec \\ \small Encadrant : Arthur Tenenhaus}
\date[juin 2026]{juin 2026}

\AtBeginSection[]{
  \begin{frame}{Plan}
    \tableofcontents[currentsection]
  \end{frame}
}

\begin{document}

% Slide 1
\begin{frame}[plain]
  \titlepage
\end{frame}

% Slide 2
\begin{frame}{Plan de la présentation}
  \tableofcontents
\end{frame}

% Slide 3: Section 1
\section{Introduction et problématique}

% Slide 4
\begin{frame}{Contexte clinique et problématique}
  \begin{block}{Le problème}
    Prédire la \textbf{localisation anatomique} d’un gliome pédiatrique à partir de la seule \textbf{signature moléculaire}\\
    de la tumeur. La localisation conditionne fortement la résécabilité chirurgicale, le protocole de\\
    radiothérapie et le pronostic.
  \end{block}

  \bigskip
  Trois localisations, donc trois classes :
  \begin{itemize}
    \item \textbf{cort} : tumeurs corticales / hémisphériques ;
    \item \textbf{dipg} : gliome pontique intrinsèque diffus, situé dans le pont de Varole (inopérable, pronostic très défavorable) ;
    \item \textbf{midl} : tumeurs de la ligne médiane, hors du pont (\textbf{classe rare}).
  \end{itemize}
\end{frame}

% Slide 5
\begin{frame}{Le jeu de données}
  \begin{columns}[T]
    \begin{column}{0.48\textwidth}
      \textbf{Deux blocs moléculaires :}\\
      \begin{itemize}
        \item \textcolor{geblue}{\textbf{GE}} : expression génique : \textbf{15 702} variables\\
              (signal discriminant principal) ;
        \item \textcolor{cghred}{\textbf{CGH}} : gains / pertes chromosomiques :\\
              \textbf{1 229} variables (signal faible isolément).
      \end{itemize}

      \medskip
      \textbf{Effectifs :} $n = 53$ patients, soit \textbf{39} en entraînement et \textbf{14} en test.

      \medskip
      \begin{center}
        \small
        \begin{tabular}{lccc}
          \toprule
                         & \textbf{cort} & \textbf{dipg} & \textbf{midl}         \\
          \midrule
          \textbf{Train} & 15            & 16            & \textcolor{geblue}{8} \\
          \textbf{Test}  & 5             & 6             & \textcolor{geblue}{3} \\
          \bottomrule
        \end{tabular}
      \end{center}
    \end{column}

    \begin{column}{0.50\textwidth}
      \centering
      \begin{tikzpicture}
        \begin{axis}[
            ybar,
            width=7.3cm,
            height=6.3cm,
            bar width=15pt,
            ymin=0, ymax=20,
            ytick={0, 5, 10, 15, 20},
            ylabel={Nb patients},
            symbolic x coords={cort,dipg,midl},
            xtick=data,
            nodes near coords,
            nodes near coords style={/pgf/number format/fixed, /pgf/number format/precision=2},
            legend style={legend columns=2},
            enlarge x limits=0.25,
            tick label style={font=\footnotesize}
          ]
          \addplot[fill=geblue] coordinates {(cort, 15) (dipg, 16) (midl, 8)};
          \addplot[fill=lightgeblue] coordinates {(cort, 5) (dipg, 6) (midl, 3)};
          \legend{Train, Test}
        \end{axis}
      \end{tikzpicture}
      \vskip 2pt
        {\footnotesize Répartition des classes : midl nettement minoritaire.}
    \end{column}
  \end{columns}
\end{frame}

% Slide 6
\begin{frame}{Le défi statistique : trois verrous}
  \begin{columns}[T]
    \begin{column}{0.58\textwidth}
      \begin{enumerate}
        \item \textbf{Régime $p \gg n$ extrême} : $p = 16\,931$ variables pour $n = 39$ patients, soit $\frac{p}{n} \simeq 434$. La matrice de Gram $X^\top X$ est de rang $\le n - 1 = 38$, donc \emph{massivement singulière} : tout estimateur non régularisé est non identifiable.
        \item \textbf{Deux blocs hétérogènes et déséquilibrés} : GE concentre $92{,}7\,\%$ des variables ; la fusion doit éviter de marginaliser CGH.
        \item \textbf{Classe rare midl} : 8 patients en train, 3 en test. Chaque erreur déplace son rappel par pas de $1/3 \simeq 0{,}33$.
      \end{enumerate}
    \end{column}

    \begin{column}{0.40\textwidth}
      \begin{block}{Réponses méthodologiques}
        \begin{itemize}
          \item Régularisation $\ell_1 /\ell_2$ + sélection parcimonieuse ;
          \item mise à l’échelle par bloc + projection supervisée (SGCCA) ;
          \item pondération de la vraisemblance + \textbf{balanced accuracy}.
        \end{itemize}
      \end{block}
    \end{column}
  \end{columns}

  \vfill
  \begin{center}
    \begin{tikzpicture}[
      box/.style={rectangle, rounded corners=3pt, draw=black, align=center, font=\scriptsize, inner sep=4pt},
      arr/.style={-{Stealth[length=2.5mm]}, thick}
      ]
      \node[box, fill=boxblue] (ge) {Bloc GE \\ 15 702};
      \node[box, fill=boxgreen, right=1.2cm of ge, yshift=-11pt] (model) {Modèle multi-blocs \\ parcimonieux};
      \node[box, fill=boxlavender, right=1.2cm of model] (loc) {Localisation \\ \texttt{cort} / \texttt{dipg} / \texttt{midl}};
      \node[box, fill=boxred, below=6pt of ge] (cgh) {Bloc CGH \\ 1 229};

      \draw[arr] (ge.east) -- (model.west);
      \draw[arr] (cgh.east) -- (model.west);
      \draw[arr] (model.east) -- (loc.west);
    \end{tikzpicture}
  \end{center}
\end{frame}

% Slide 7: Section 2
\section{Analyse exploratoire (EDA)}

% Slide 8
\begin{frame}{Domination du bloc GE et pondération des classes}
  \vspace{14pt}
  \begin{columns}[c]
    \begin{column}{0.48\textwidth}
      \centering
      \hspace{5.5pt}%
      \begin{tikzpicture}
        \begin{axis}[
            ybar,
            width=5.22cm,
            height=3.40cm,
            scale only axis,
            bar width=23.9pt,
            ymin=0, ymax=108,
            ytick={0, 50, 100},
            ylabel={\% variables},
            symbolic x coords={GE,CGH},
            xtick={GE,CGH},
            nodes near coords,
            nodes near coords style={/pgf/number format/fixed, /pgf/number format/precision=1},
            enlarge x limits={abs=40.34pt},
          ]
          \addplot[fill=geblue, bar shift=0pt] coordinates {(GE, 92.7)};
          \addplot[fill=cghred, bar shift=0pt] coordinates {(CGH, 7.3)};
        \end{axis}
      \end{tikzpicture}
      \vskip 2pt
        {\footnotesize Inertie $\propto$ nombre de variables : la concaténation brute est dominée à > 90\% par GE.}
    \end{column}

    \begin{column}{0.50\textwidth}
      \centering
      \begin{tikzpicture}
        \begin{axis}[
            ybar,
            width=5.22cm,
            height=3.40cm,
            scale only axis,
            bar width=21.9pt,
            ymin=0, ymax=2,
            ytick={0, 0.5, 1, 1.5, 2},
            ylabel={poids $w_c$},
            symbolic x coords={cort,dipg,midl},
            xtick={cort,dipg,midl},
            nodes near coords,
            nodes near coords style={/pgf/number format/fixed, /pgf/number format/precision=2},
            enlarge x limits={abs=27.74pt},
          ]
          \addplot[fill=barlightblue, bar shift=0pt] coordinates {(cort, 0.87) (dipg, 0.81)};
          \addplot[fill=cghred, bar shift=0pt] coordinates {(midl, 1.62)};
        \end{axis}
      \end{tikzpicture}
      \vskip 2pt
        {\footnotesize Poids inverses de prévalence $w_c = n/(K\,n_c)$ : un patient\\ midl pèse $\simeq 1{,}6\times$ un patient moyen.}
    \end{column}
  \end{columns}
\end{frame}

% Slide 9
\begin{frame}{Qualité des données : niveaux d’expression GE déjà normalisés}
  \begin{columns}[T]
    \begin{column}{0.48\textwidth}
      \vskip 5pt
      \centering
      \begin{tikzpicture}
        \begin{axis}[
            ybar,
            width=5.80cm,
            height=3.41cm,
            scale only axis,
            bar width=4.98pt,
            ymin=0, ymax=72000,
            ytick={0, 20000, 40000, 60000},
            yticklabels={0, 20k, 40k, 60k},
            xmin=-11, xmax=11,
            xtick={-10, -5, 0, 5, 10},
            xlabel={Niveau d’expression normalisé},
            ylabel={Fréquence},
            scaled y ticks=false,
          ]
          \addplot[fill=geblue, draw=geblue, bar shift=0pt] coordinates {
              (-10, 80) (-9, 160) (-8, 320) (-7, 650) (-6, 1300) (-5, 2600)
              (-4, 5200) (-3, 11000) (-2, 24000) (-1, 46000) (0, 64000)
              (1, 45000) (2, 23000) (3, 10500) (4, 5000) (5, 2500)
              (6, 1250) (7, 620) (8, 300) (9, 150) (10, 70)
            };
        \end{axis}
      \end{tikzpicture}
    \end{column}

    \begin{column}{0.50\textwidth}
      Distribution globale des $39 \times 15\,702$ valeurs d’expression :
      \begin{itemize}
        \item quasi symétrique, centrée (médiane = 0,03), bornée ($-10{,}43$ à $9{,}87$) ;
        \item données déjà normalisées (log-ratios) $\Rightarrow$ aucune transformation (log, Box-Cox), une simple standardisation z-score suffit.
      \end{itemize}
      \medskip
      \textit{Imputation} par la médiane du train (reportée au test) : robuste aux valeurs extrêmes, sans fuite d’information.
    \end{column}
  \end{columns}
\end{frame}

% Slide 10
\begin{frame}{La PCA non supervisée ne sépare pas les classes}
  \begin{columns}[c]
    \begin{column}{0.48\textwidth}
      \centering
      \begin{tikzpicture}
        \begin{axis}[
            scale only axis,
            width=4.6cm,
            height=3.0cm,
            xmin=-1.5, xmax=1.5,
            ymin=-1.5, ymax=1.5,
            xtick={-1, 0, 1},
            ytick={-1, 0, 1},
            xlabel={\footnotesize PC1 (10,8\%)},
            ylabel={\footnotesize PC2 (9,4\%)},
            legend pos=outer north east,
            legend style={font=\tiny},
            tick label style={font=\tiny}
          ]
          \addplot[only marks, mark=*, mark size=1.5pt, geblue] coordinates {
              (-1.17, 0.60) (-0.47, 1.10) (0.33, -0.40) (1.03, 0.80) (0.63, -1.00) (-0.77, -0.30) (1.43, 0.20) (0.13, 1.30)
            };
          \addplot[only marks, mark=*, mark size=1.5pt, dipgorange] coordinates {
              (0.23, 0.30) (0.93, -0.60) (-0.27, 0.90) (1.13, 1.00) (-0.67, 0.40) (0.53, -1.10) (1.33, -0.20) (-0.97, 0.70)
            };
          \addplot[only marks, mark=*, mark size=1.5pt, midlbrown] coordinates {
              (0.04, 0.10) (0.44, 0.50) (-0.16, -0.30) (0.74, -0.20) (-0.36, 0.60)
            };
          \legend{cort, dipg, midl}
        \end{axis}
      \end{tikzpicture}
    \end{column}

    \begin{column}{0.48\textwidth}
      \begin{itemize}
        \item PC1+PC2 ne capturent que $\simeq 20\%$ de l’inertie pour GE ($\simeq 27\%$ pour CGH) ;
        \item les trois nuages se \textbf{chevauchent fortement} ; \texttt{midl} ne forme \textit{aucun amas distinct}.
      \end{itemize}
      \medskip
      \begin{alertblock}{Conséquence}
        Une réduction \textit{non supervisée} (PCA, Sparse PCA) jette la variance dans des directions non discriminantes $\Rightarrow$ échec des pipelines Sparse PCA + LogReg. \textbf{Il faut une projection supervisée} (SGCCA).
      \end{alertblock}
    \end{column}
  \end{columns}
\end{frame}

% Slide 11
\begin{frame}{Signal univarié réel mais peu robuste, et forte redondance}
  \begin{columns}[T]
    \begin{column}{0.50\textwidth}
      \boldmath{$t$} \textbf{de Welch univariés} \texttt{\textbf{midl}} \textbf{vs reste} (Welch car variances inégales, $n_{\text{midl}} = 8$) :
      \begin{itemize}
        \item \textbf{GE} : quelques gènes à $|t| \simeq 4$–$5$ (p.ex. \texttt{A\_32\_P85999}) $\Rightarrow$ un signal univarié existe ;
        \item \textbf{CGH} : statistiques plafonnant vers $|t| \simeq 3$, signal plus faible.
      \end{itemize}
      \bigskip
      \textbf{Forte redondance intra-bloc} (groupes de variables corrélées, surtout CGH) :
      \begin{itemize}
        \item justifie la pénalisation \textbf{Elastic-Net}
      \end{itemize}
    \end{column}

    \begin{column}{0.48\textwidth}
      \begin{exampleblock}{Univarié $\neq$ multivarié robuste}
        Avec $p \gg n$ et $n_{\text{midl}} = 8$, ces t-statistiques sont elles-mêmes très bruitées (un patient peut faire basculer le classement) et ne contrôlent \textbf{pas} la multiplicité ($\sim 17\,000$ tests). Diagnostic exploratoire, pas une sélection valide $\Rightarrow$ sélection confiée à des méthodes pénalisées \textit{intra-CV}.
      \end{exampleblock}
    \end{column}
  \end{columns}
\end{frame}

% Slide 12
\begin{frame}{Signal univarié autour de \texttt{midl} : top variables ($t$ de Welch)}
  \begin{columns}[T]
    \begin{column}{0.56\textwidth}
      \centering
      \begin{tikzpicture}
        \begin{axis}[
            xbar,
            bar shift=0pt,
            scale only axis,
            width=159.26pt,
            height=108.24pt,
            bar width=7pt,
            xmin=-5.5, xmax=5.5,
            xtick={-4, -2, 0, 2, 4},
            xlabel={$t$ de Welch (\texttt{midl} vs reste)},
            ymin=0.4, ymax=8.6,
            ytick={1, 2, 3, 4, 5, 6, 7, 8},
            yticklabels={{A\_23\_P37560}, {A\_23\_P163278}, {A\_23\_P37475}, {A\_32\_P430359}, {A\_24\_P349151}, {A\_23\_P129188}, {A\_24\_P936875}, {A\_32\_P85999}},
            yticklabel style={font=\tiny\ttfamily},
            nodes near coords,
            nodes near coords style={/pgf/number format/fixed, font=\tiny},
            tick label style={font=\tiny}
          ]
          \addplot[draw=geblue, fill=geblue] coordinates {
              (-3.4, 1)
              (-3.6, 2)
              (-3.7, 3)
              (-3.9, 4)
              (-4, 5)
              (-4.1, 6)
            };
          \addplot[draw=midlbrown, fill=midlbrown] coordinates {
              (4.1, 7)
              (4.4, 8)
            };
        \end{axis}
      \end{tikzpicture}
    \end{column}

    \begin{column}{0.42\textwidth}
      \footnotesize
      \begin{itemize}
        \item \textbf{GE} : quelques gènes atteignent $|t| \simeq 4$–$5$ ; rouge = sur-exprimé chez \texttt{midl} (p.ex. \texttt{A\_32\_P85999}), la majorité y étant sous-exprimés ;
        \item \textbf{CGH} : statistiques plafonnant vers $|t| \simeq 3$, cohérent avec sa quasi-exclusion par les modèles parcimonieux.
      \end{itemize}
    \end{column}
  \end{columns}
\end{frame}

% Slide 13
\begin{frame}{Choix de la métrique : balanced accuracy}
  \textit{L’accuracy} brute pondère implicitement par la fréquence et masquerait un échec sur \texttt{midl} (un classifieur ne prédisant jamais midl atteint déjà $\simeq 11/14$).

  \medskip
  \begin{columns}[c]
    \begin{column}{0.54\textwidth}
      On retient le \textbf{macro-recall} (= balanced accuracy) :
      \[
        \text{BA}(\hat{f}) = \frac{1}{|\mathcal{Y}|} \sum_{c \in \mathcal{Y}} \underbrace{\frac{|\{i : Y_i = c, \hat{f}(X_i) = c\}|}{|\{i : Y_i = c\}|}}_{\text{recall}_c}
      \]
      même poids aux trois classes, indépendamment de la prévalence.
    \end{column}

    \begin{column}{0.44\textwidth}
      \begin{alertblock}{Un exemple}
        Pour un classifieur ne prédisant \textit{jamais} \texttt{midl} :
        \[
          \text{macro-recall} = \frac{1}{3}\left(\frac{5}{5} + \frac{6}{6} + \frac{0}{3}\right) = 0{,}667
        \]
      \end{alertblock}
    \end{column}
  \end{columns}
\end{frame}

% Slide 14: Section 3
\section{Formalisme et méthodes}

% Slide 15
\begin{frame}{Cadre probabiliste et classifieur de Bayes}
  \vspace{7pt}
  Classification à 3 classes : $(X_i, Y_i) \text{ i.i.d.} \sim P$, $X \in \mathbb{R}^p$, $Y \in \{1, 2, 3\}$. Risque d’un classifieur $f$ pour une perte $\ell$ :
  \[
    R(f) = \mathbb{E}_{(X,Y)\sim P}[\ell(f(X), Y)], \qquad f^\star = \arg\min_f R(f).
  \]

  \bigskip
  \begin{block}{Sous la perte 0--1}
    Avec $\eta_c(x) = \mathbb{P}[Y = c \mid X = x]$, le prédicteur de Bayes est
    \[
      f^\star(x) = \arg\max_{c \in \mathcal{Y}} \eta_c(x).
    \]
    La 0--1 loss n’étant ni continue ni dérivable, on utilise une \textbf{perte de substitution} convexe et différentiable : l’entropie croisée.
  \end{block}

  \bigskip
  \textit{Remarque :} sous une perte \textit{asymétrique} (pénaliser davantage l’erreur sur \texttt{midl}), le prédicteur de Bayes n’est plus l’argmax \\
  des $\eta_c$ : ce qui motive la pondération.
\end{frame}

% Slide 16
\begin{frame}{Régression logistique multinomiale}
  Modèle softmax (avec contrainte d’identifiabilité $\beta_3 = 0$) :
  \[
    \eta_c^\beta(x) = \mathbb{P}_\beta[Y = c \mid X = x] = \frac{e^{\beta_c^\top x}}{\sum_{k=1}^3 e^{\beta_k^\top x}}.
  \]
  Maximiser la log-vraisemblance $\equiv$ minimiser l’\textbf{entropie croisée} empirique :
  \[
    \hat{R}_n(\beta) = -\frac{1}{n} \sum_{i=1}^n \log \eta_{Y_i}^\beta(X_i), \qquad \nabla_{\beta_c} \hat{R}_n(\beta) = \frac{1}{n} \sum_{i=1}^n \left( \eta_c^\beta(X_i) - \mathbf{1}_{\{Y_i = c\}} \right) X_i.
  \]
  \begin{itemize}
    \item Objectif \textbf{convexe et différentiable} ; pas de solution analytique $\Rightarrow$ Méthode de Newton.
    \item Frontière de décision \textbf{linéaire} : $\hat{y} = \arg\max_c \hat{\beta}_c^\top x$.
  \end{itemize}
\end{frame}

% Slide 17
\begin{frame}{Vraisemblance pondérée : poids inverses de prévalence}
  Le risque uniforme $\hat{R}_n$ est dominé par les classes majoritaires. On aligne l’objectif sur la balanced accuracy via
  \[
    \hat{R}_n^w(\beta) = -\frac{1}{n}\sum_{i=1}^n w_{Y_i} \log \eta_{Y_i}^\beta(X_i).
  \]

  \vspace{5pt}
  \begin{block}{Poids}
    Égaliser la contribution de chaque classe ($w_c n_c = \text{cte}$) puis normaliser ($\sum_c w_c n_c = n$) donne
    \[
      w_c = \frac{n}{K\,n_c}, \qquad w_{\texttt{cort}} = 0{,}87, \quad w_{\texttt{dipg}} = 0{,}81, \quad w_{\texttt{midl}} = 1{,}62.
    \]
  \end{block}

  \bigskip
  Le gradient pondéré amplifie d’un facteur $w_{\texttt{midl}}$ les résidus \texttt{midl}. Plus défendable qu’un sur-échantillonnage \texttt{SMOTE} (testé, peu concluant : fabriquer de faux patients en dimension 16 931 est instable). Le problème reste \textbf{convexe}.
\end{frame}

% Slide 18
\begin{frame}{Pénalisation Elastic-Net et cadre GLM}
  \begin{columns}[T]
    \begin{column}{0.50\textwidth}
      \textbf{Elastic-Net} pour contrôler la complexité en $p \gg n$ :
      \[
        \min_\beta \left\{ \hat{R}_n(\beta) + \lambda\rho\|\beta\|_1 + \frac{\lambda(1-\rho)}{2}\|\beta\|_2^2 \right\}
      \]
      \vspace{0.2cm}
      \begin{itemize}
        \item $\ell_1$ (Lasso) $\Rightarrow$ \textbf{sélection} de variables
        \item $\ell_2$ (Ridge) $\Rightarrow$ \textbf{stabilité} face aux variables colinéaires
        \item $\ell_1$ non différentiable $\Rightarrow$ optimisation proximale (soft-thresholding).
      \end{itemize}
    \end{column}

    \begin{column}{0.48\textwidth}
      \begin{block}{GLM}
        La régression logistique est le \textit{GLM} canonique : $Y \mid X \sim \text{loi catégorielle}$, lien logit, $g(\mu_i) = x_i^\top \beta$. La log-vraisemblance est concave $\Rightarrow$ l'EMV est unique et s'obtient par Newton (moindres carrés repondérés, IRLS).
      \end{block}
    \end{column}
  \end{columns}
\end{frame}

% Slide 19
\begin{frame}{Analyse Discriminante Linéaire (LDA)}
  Hypothèses : normalité conditionnelle $X \mid Y = c \sim \mathcal{N}_p(\mu_c, \Sigma)$ et \textbf{homoscédasticité} ($\Sigma_c = \Sigma$). La fonction discriminante se linéarise :
  \[
    \delta_c(x) = \underbrace{(\Sigma^{-1}\mu_c)^\top}_{w_c^\top} x - \underbrace{\left(\tfrac{1}{2}\mu_c^\top \Sigma^{-1}\mu_c - \log \pi_c\right)}_{b_c}, \qquad f^\star(x) = \arg\max_c w_c^\top x - b_c.
  \]
  Estimateurs : $\hat{\pi}_c = \frac{n_c}{n}$, $\hat{\mu}_c$ moyenne de classe, et covariance commune
  \[
    \hat{\Sigma} = \frac{1}{n-K} \sum_c \sum_i (X_i - \hat{\mu}_c)(X_i - \hat{\mu}_c)^\top \mathbf{1}_{Y_i = c}.
  \]

  \medskip
  \begin{alertblock}{Problème en $p \gg n$}
    $\hat{\Sigma}$ ($16\,931 \times 16\,931$) est \textbf{non inversible}. La LDA n'est donc pas appliquée sur les variables brutes mais sur les \textit{composantes supervisées} issues de la SGCCA (quelques dimensions seulement).
  \end{alertblock}
\end{frame}

% Slide 20
\begin{frame}{Validation croisée : stratifiée, répétée, imbriquée}
  Avec $n = 39$, tout split unique est trop instable.

  \bigskip
  \begin{columns}[T]
    \begin{column}{0.46\textwidth}
      \begin{itemize}
        \item \textbf{Stratification} : chaque pli préserve la prévalence des classes (sinon un pli peut ne contenir aucun \texttt{midl}) ;
        \item \textbf{Répétition} : $K \times R = 7 \times 3 = 21$ plis pour réduire la variance de $\hat{R}_{\text{CV}}$ ;
        \item \textbf{Métrique} : balanced accuracy optimisée et rapportée.
      \end{itemize}
    \end{column}

    \begin{column}{0.50\textwidth}
      \begin{block}{CV imbriquée (nested)}
        \begin{itemize}
          \item \textit{Boucle externe} : évalue $\hat{R}_{\text{CV}}$ ;
          \item \textit{Boucle interne} : sélectionne $\hat{\lambda} = \arg\min_\lambda \hat{R}_{\text{int}}(\lambda)$ sur le seul fold-train.
        \end{itemize}
        $\Rightarrow \hat{\lambda}$ ne voit jamais les données d’évaluation : pas de \textbf{data leakage}. Toute sélection de variables se fait \textit{à l’intérieur} des plis.
      \end{block}
    \end{column}
  \end{columns}
\end{frame}

% Slide 21
\begin{frame}{Fusion multibloc : des méthodes naïves au Cooperative Learning}
  Deux blocs : $X \in \mathbb{R}^{n \times p_{\text{GE}}}$ et $Z \in \mathbb{R}^{n \times p_{\text{CGH}}}$.
  \begin{itemize}
    \item \textbf{Early fusion} : concaténer $[X \mid Z] \Rightarrow$ GE domine, CGH marginalisé ;
    \item \textbf{Late fusion} : un modèle par bloc puis $\omega f_X + (1-\omega) f_Z \Rightarrow$ ignore les interactions inter-blocs.
  \end{itemize}

  \begin{block}{Cooperative Learning (Ding \& Tibshirani, PNAS 2022)}
    Interpolation continue entre les deux, via une \textbf{pénalité d’accord} $\rho \ge 0$ sur les logits :
    \[
      J(\theta_x, \theta_z) = \hat{R}_{\text{CE}}(\theta_x, \theta_z) + \underbrace{\frac{\rho}{2} \|X \theta_x - Z \theta_z\|_F^2}_{\text{accord}} + P_{\text{Elastic-Net}}(\theta_x, \theta_z).
    \]
    $\rho = 0$ : early fusion ; $\rho \to \infty$ : late fusion. Problème équivalent à un problème Elastic-Net sur les données \textit{augmentées} $(\tilde{X}, \tilde{z})$.
  \end{block}
\end{frame}

% Slide 22
\begin{frame}{SGCCA : la projection supervisée retenue}
  \textbf{Sparse Generalized CCA} (Tenenhaus \textit{et al.}) : composantes parcimonieuses $y_j = X_j a_j$ maximisant la covariance \textit{entre blocs connectés} :
  \[
    \max_{a_1, \dots, a_J} \sum_{j,k} c_{jk} g\left(\text{cov}(X_j a_j, X_k a_k)\right) \quad \text{s.c.} \quad \|a_j\|_2 \le 1, \; \|a_j\|_1 \le s_j.
  \]

  \begin{columns}[c]
    \begin{column}{0.54\textwidth}
      \begin{itemize}
        \item \textbf{3 blocs} : GE, CGH, et le bloc des \textit{labels} $Y$ comme \textbf{pivot supervisé} ;
        \item matrice de connexion : on corrèle chaque bloc moléculaire à $Y$, pas GE–CGH directement
              \[
                C = \begin{pmatrix} 0 & 0 & 1 \\ 0 & 0 & 1 \\ 1 & 1 & 0 \end{pmatrix} ;
              \]
        \item budget $\ell_1 \Rightarrow$ sélection ; boule $\ell_2 \Rightarrow$ stabilité ;
        \item mise à l’échelle $\tilde{X}_j = X_j / \sqrt{p_j}$ pour égaliser l’inertie des blocs.
      \end{itemize}
    \end{column}

    \begin{column}{0.44\textwidth}
      \centering
      \begin{tikzpicture}[
        box/.style={rectangle, rounded corners=3pt, draw=black, align=center, font=\scriptsize, inner sep=3pt},
        arr/.style={-{Stealth[length=2.5mm]}, thick}
        ]
        \node[box, fill=boxblue] (ge) {GE};
        \node[box, fill=boxred, below=12pt of ge] (cgh) {CGH};
        \path (ge) -- (cgh) coordinate[midway] (mid);
        \node[box, fill=boxgreen, right=1.0cm of mid] (sgcca) {SGCCA \\ $\xi_{\text{GE}}, \xi_{\text{CGH}}$};
        \node[box, fill=boxlavender, right=0.8cm of sgcca] (lda) {LDA};
        \node[right=0.6cm of lda, font=\scriptsize] (classe) {classe};

        \draw[arr] (ge.east) -- (sgcca.west);
        \draw[arr] (cgh.east) -- (sgcca.west);
        \draw[arr] (sgcca.east) -- (lda.west);
        \draw[arr] (lda.east) -- (classe.west);
      \end{tikzpicture}
      \vskip 4pt
        {\footnotesize Les composantes alimentent une LDA dans un espace réduit, contournant la non-inversibilité de $\hat{\Sigma}$.}
    \end{column}
  \end{columns}
\end{frame}

% Slide 23: Section 4
\section{Résultats}

% Slide 24
\begin{frame}{Sélection de variables et comparaison des pipelines}
  Quatre pipelines sparses supervisés, même protocole (CV stratifiée $7 \times 3 = 21$ plis, balanced accuracy, test $n = 14$ tenu à part). SGCCA retient \textbf{68 GE + 11 CGH} ($s_1 = 0{,}051$, $s_2 = 0{,}067$).

  \medskip
  \begin{center}
    \small
    \begin{tabular}{lcccc}
      \toprule
      \textbf{Méthode}                     & \textbf{CV bal\_acc} & \textbf{Test bal\_acc} & \texttt{\textbf{midl}} \textbf{test} & \textbf{N vars} \\
      \midrule
      Multinomial Lasso \textit{grouped}   & 0,784 $\pm$ 0,096    & 0,773                  & 0/3                                  & 27 GE           \\
      Multinomial Lasso \textit{ungrouped} & 0,838 $\pm$ 0,123    & 0,773                  & 0/3                                  & 27 GE           \\
      \rowcolor{boxgreen}
      Cooperative Lasso OvR ($\rho = 0$)   & 0,833 $\pm$ 0,129    & 0,924                  & 2/3                                  & 42 GE           \\
      \rowcolor{boxgreen}
      SGCCA + LDA                          & 0,829 $\pm$ 0,133    & 0,924                  & 2/3                                  & 68 GE + 11 CGH  \\
      \bottomrule
    \end{tabular}
  \end{center}

  \bigskip
  \begin{itemize}
    \item Seuls les pipelines \textit{One-vs-Rest} ou à \textit{espace latent supervisé} récupèrent \texttt{midl} (2/3) ;
    \item passer la pénalité $\ell_1$ \textit{grouped} $\to$ \textit{ungrouped} fait gagner +0,054 en CV (au seul prix d’un argument) mais rien sur le test.
  \end{itemize}
\end{frame}

% Slide 25
\begin{frame}{Performance comparée}
  \begin{columns}[c]
    \begin{column}{0.48\textwidth}
      \centering
      {\footnotesize Balanced accuracy en CV ($\pm$ écart-type)}
      \begin{tikzpicture}
        \begin{axis}[
            ybar,
            width=5.5cm,
            height=4.2cm,
            bar width=10pt,
            ymin=0.3, ymax=1.05,
            ytick={0.4, 0.6, 0.8, 1},
            ylabel={\footnotesize CV bal\_acc},
            symbolic x coords={SGCCA,Coop,M-grp,M-ungrp},
            xtick=data,
            nodes near coords,
            nodes near coords style={/pgf/number format/fixed, /pgf/number format/precision=2, font=\tiny},
            enlarge x limits=0.3,
            tick label style={font=\tiny}
          ]
          \addplot[
            fill=midlgreen,
            draw=black,
            error bars/.cd,
            y dir=both,
            y explicit
          ] coordinates {
              (SGCCA, 0.829) +- (0, 0.133)
              (Coop, 0.833) +- (0, 0.129)
              (M-grp, 0.784) +- (0, 0.096)
              (M-ungrp, 0.838) +- (0, 0.123)
            };
        \end{axis}
      \end{tikzpicture}
    \end{column}

    \begin{column}{0.48\textwidth}
      \centering
      {\footnotesize Balanced accuracy sur le test ($n = 14$)}
      \begin{tikzpicture}
        \begin{axis}[
            ybar,
            width=5.5cm,
            height=4.2cm,
            bar width=10pt,
            ymin=0.3, ymax=1.05,
            ytick={0.4, 0.6, 0.8, 1},
            ylabel={\footnotesize Test bal\_acc},
            symbolic x coords={SGCCA,Coop,M-grp,M-ungrp},
            xtick=data,
            nodes near coords,
            nodes near coords style={/pgf/number format/fixed, /pgf/number format/precision=2, font=\tiny},
            enlarge x limits=0.3,
            tick label style={font=\tiny}
          ]
          \addplot[fill=geblue, draw=black] coordinates {
              (SGCCA, 0.924)
              (Coop, 0.924)
              (M-grp, 0.773)
              (M-ungrp, 0.773)
            };
        \end{axis}
      \end{tikzpicture}
    \end{column}
  \end{columns}

  \vfill
  \begin{center}
    \footnotesize En CV, trois méthodes plafonnent à $\simeq 0{,}83$ ; sur le test, seules \textbf{SGCCA + LDA} et \textbf{Cooperative OvR} atteignent 0,924.
  \end{center}
\end{frame}

% Slide 26
\begin{frame}{Parcimonie effective : combien de variables par bloc ?}
  \begin{columns}[T]
    \begin{column}{0.54\textwidth}
      \centering
      \begin{tikzpicture}
        \begin{axis}[
            ybar=2pt,
            width=5.8cm,
            height=3.8cm,
            scale only axis,
            bar width=11pt,
            ymin=0, ymax=80,
            ytick={0, 20, 40, 60, 80},
            ylabel={\footnotesize Nb variables retenues},
            symbolic x coords={SGCCA,Coop,M-grp,M-ungrp},
            xtick=data,
            nodes near coords,
            nodes near coords style={/pgf/number format/fixed, font=\tiny},
            legend style={at={(0.5, 1.03)}, anchor=south, legend columns=2, font=\tiny},
            enlarge x limits=0.3,
            tick label style={font=\tiny}
          ]
          \addplot[fill=geblue] coordinates {(SGCCA, 68) (Coop, 42) (M-grp, 27) (M-ungrp, 27)};
          \addplot[fill=cghred] coordinates {(SGCCA, 11) (Coop, 0) (M-grp, 0) (M-ungrp, 0)};
          \legend{GE (sur 15 702), CGH (sur 1 229)}
        \end{axis}
      \end{tikzpicture}
    \end{column}

    \begin{column}{0.44\textwidth}
      Toutes les méthodes \textit{data-driven} excluent spontanément l’intégralité du bloc CGH.

      \medskip
      Seule la \textbf{SGCCA} en conserve 11 variables : et uniquement parce que la contrainte de parcimonie $s_{\text{CGH}}$ est imposée structurellement bloc par bloc.

      \medskip
      Modèles très parcimonieux et interprétables ($\le 68$ variables sur $\sim 17\,000$).
    \end{column}
  \end{columns}
\end{frame}

% Slide 27
\begin{frame}{Le Cooperative Learning : une instabilité instructive}
  Avec \texttt{multiview} (sur \texttt{glmnet}), l’algorithme \textbf{diverge pour tout} $\rho > 0$ en formulation OvR (\texttt{algorithm did not converge}). Deux causes :
  \begin{itemize}
    \item les poids de Newton $w_i = \hat{p}_i(1 - \hat{p}_i) \to 0$ sur la classe rare $\Rightarrow X^\top W X$ quasi-singulière ;
    \item l’accord ajoute $\rho X^\top X$ et amplifie d’un facteur $(1 + \rho)$ des corrélations intra-bloc déjà extrêmes.
  \end{itemize}

  \bigskip
  \begin{block}{Lecture structurelle}
    La pénalité $\ell_1$ pousse à \textit{annuler} des coefficients tandis que l’accord pousse à \textit{coupler} les vues : forces contradictoires quand GE et CGH ne partagent pas de signal commun fort. Seul $\rho = 0$ est exploitable (la méthode dégénère en Lasso OvR). La formulation \textit{multinomiale} converge ($\rho^\star = 0{,}10$) mais à balanced accuracy plus basse.
  \end{block}
\end{frame}

% Slide 28
\begin{frame}{Sensibilité aux hyperparamètres}
  \begin{columns}[T]
    \begin{column}{0.48\textwidth}
      \centering
      {\footnotesize \textbf{Multinomial Lasso : sweep de $\alpha$}}
      \begin{tikzpicture}
        \begin{axis}[
            width=5.5cm,
            height=4.0cm,
            xmin=0, xmax=1,
            ymin=0.6, ymax=0.9,
            xtick={0, 0.5, 1},
            ytick={0.6, 0.7, 0.8, 0.9},
            xlabel={\footnotesize $\alpha$ (Lasso–Ridge)},
            ylabel={\footnotesize CV bal\_acc},
            legend style={at={(0.7,0.95)}, anchor=north, legend columns=2, font=\tiny},
            tick label style={font=\tiny}
          ]
          \addplot[thick, gebluemid, mark=*] coordinates {(0, 0.706) (0.5, 0.781) (1.0, 0.784)};
          \addplot[thick, dipgorange, mark=*] coordinates {(0, 0.706) (0.5, 0.794) (1.0, 0.838)};
          \legend{grouped, ungrouped}
        \end{axis}
      \end{tikzpicture}
      \vskip 2pt
        {\footnotesize À $\alpha = 1$, passer \textit{grouped} $\to$ \textit{ungrouped} : +0,054.}
    \end{column}

    \begin{column}{0.48\textwidth}
      \centering
      {\footnotesize \textbf{Cooperative OvR : sweep de $\rho$}}
      \begin{tikzpicture}
        \begin{axis}[
            width=5.5cm,
            height=4.0cm,
            xmin=0, xmax=5,
            ymin=0, ymax=0.9,
            xtick={0, 1, 2, 3, 4, 5},
            ytick={0, 0.4, 0.6, 0.8},
            xlabel={\footnotesize $\rho$ (accord)},
            ylabel={\footnotesize CV bal\_acc},
            tick label style={font=\tiny}
          ]
          \addplot[thick, gebluemid, mark=*] coordinates {(0, 0.833)};
          \addplot[only marks, mark=x, mark size=3pt, red, thick] coordinates {
              (1, 0.4) (2, 0.4) (3, 0.4) (4, 0.4) (5, 0.4)
            };
        \end{axis}
      \end{tikzpicture}
      \vskip 2pt
        {\footnotesize Croix rouges : non-convergence du solveur (\texttt{multiview}/\texttt{glmnet}) ;\\ seul $\rho = 0$ exploitable.}
    \end{column}
  \end{columns}
\end{frame}

% Slide 29
\begin{frame}{Multinomial vs OvR : le verrou de la classe rare}
  \vspace{15pt}
  Mêmes équations, comportements différents en $p \gg n$ avec classes déséquilibrées.

  \medskip
  \begin{columns}[c]
    \begin{column}{0.48\textwidth}
      \centering
      \textbf{SGCCA + LDA / Cooperative OvR}
      \vskip 2pt
      \small
      \setlength{\tabcolsep}{7pt}
      \begin{tabular}{lccc}
                        & p.\texttt{cort}          & p.\texttt{dipg}          & p.\texttt{midl}                 \\
        v.\texttt{cort} & \cellcolor{lightgeblue}5 & 0                        & 0                               \\
        v.\texttt{dipg} & 0                        & \cellcolor{lightgeblue}6 & 0                               \\
        v.\texttt{midl} & 0                        & 1                        & \cellcolor{cellgreen}\textbf{2} \\
      \end{tabular}
      \vskip 2pt
      bal\_acc = 0,924 : \texttt{midl} 2/3
    \end{column}

    \begin{column}{0.48\textwidth}
      \centering
      \textbf{Multinomial Lasso (grouped/ungrouped)}
      \vskip 2pt
      \small
      \setlength{\tabcolsep}{7pt}
      \begin{tabular}{lccc}
                        & p.\texttt{cort}          & p.\texttt{dipg}          & p.\texttt{midl}           \\
        v.\texttt{cort} & \cellcolor{lightgeblue}5 & 0                        & 0                         \\
        v.\texttt{dipg} & 0                        & \cellcolor{lightgeblue}6 & 0                         \\
        v.\texttt{midl} & 1                        & \cellcolor{cellred}2     & \cellcolor{celllightred}0 \\
      \end{tabular}
      \vskip 2pt
      bal\_acc = 0,773 : \texttt{midl} 0/3
    \end{column}
  \end{columns}

  \medskip
  \begin{exampleblock}{Le « filet de sauvetage » de l’OvR}
    En OvR, un classifieur \texttt{midl}-vs-reste à coefficients nuls garde un \textit{intercept} calibré sur la prévalence :\\
    $\hat{P}_{\texttt{midl}} \simeq \sigma(\log(8/31)) \simeq 0{,}205$ \textit{constant}.
    Quand \texttt{cort} et \texttt{dipg} rejettent un patient, ce plancher l’emporte par $\arg\max$ (p.ex. P11, P19). La contrainte $\sum_c P_c = 1$ du multinomial aspire au contraire cette masse vers \texttt{dipg}.
  \end{exampleblock}
\end{frame}

% Slide 30
\begin{frame}{Repêchage de midl : probabilités prédites (Cooperative OvR)}
  \begin{columns}[T]
    \begin{column}{0.54\textwidth}
      \centering
      \begin{tikzpicture}
        \begin{axis}[
            ybar=2pt,
            width=5.8cm,
            height=3.6cm,
            scale only axis,
            bar width=8pt,
            ymin=0, ymax=0.3,
            ytick={0, 0.1, 0.2, 0.3},
            ylabel={\footnotesize $\hat{P}(\text{classe} \mid x)$},
            symbolic x coords={P11,P14,P19},
            xtick=data,
            nodes near coords,
            nodes near coords style={/pgf/number format/fixed, /pgf/number format/precision=2, /pgf/number format/zerofill=true, font=\tiny},
            legend style={at={(0.5, 1.03)}, anchor=south, legend columns=3, font=\tiny},
            enlarge x limits=0.2,
            xticklabel style={font=\small},
            yticklabel style={font=\tiny}
          ]
          \addplot[fill=geblue] coordinates {(P11, 0.02) (P14, 0.04) (P19, 0.03)};
          \addplot[fill=dipgorange] coordinates {(P11, 0.10) (P14, 0.18) (P19, 0.08)};
          \addplot[fill=midlbrown] coordinates {(P11, 0.20) (P14, 0.20) (P19, 0.20)};
          \legend{cort, dipg, midl}
        \end{axis}
      \end{tikzpicture}
    \end{column}

    \begin{column}{0.44\textwidth}
      Pour les 3 vrais \texttt{midl} du test, l’OvR prédit un\\
      plancher quasi constant
      \[
        \hat{P}_{\texttt{midl}} \simeq \sigma(\log(8/31)) \simeq 0{,}205.
      \]
      \texttt{cort} et \texttt{dipg} étant rejetés, ce plancher l’emporte
      par \textbf{argmax} (P11, P19).

      \bigskip
      Les multinomiaux, contraints par $\sum_c P_c = 1$,\\
      aspirent cette masse vers \texttt{dipg} $\Rightarrow$ recall \texttt{midl}\\
      = 0/3.
    \end{column}
  \end{columns}
\end{frame}

% Slide 31: Section 5
\section{Stabilité et reproductibilité}

% Slide 32
\begin{frame}{Stabilité de la sélection de variables}
  En $p \gg n$, l’identité exacte des variables retenues est instable. On la quantifie par \textbf{ré-échantillonnage\\
    bootstrap} (Efron, 1979) : fréquence de sélection $\hat{\Pi}_k = \frac{1}{B} \sum_b \mathbf{1}_{\{a_k^{(b)} \ne 0\}}$, « stable set »\\
  $= \{k : \hat{\Pi}_k \ge \pi_{\text{thr}}\}$ (Meinshausen \& Bühlmann, 2010 ; $\pi_{\text{thr}} = 0{,}6$, $B = 200$).

  \bigskip
  \begin{columns}[T]
    \begin{column}{0.50\textwidth}
      \centering
      \begin{tikzpicture}
        \begin{axis}[
            xbar,
            width=5.8cm,
            height=4.0cm,
            bar width=6pt,
            xmin=0, xmax=1,
            xtick={0, 0.2, 0.4, 0.6, 0.8, 1},
            xlabel={\footnotesize $\hat{\Pi}_k$},
            ymin=0.5, ymax=4.5,
            ytick={1, 2, 3, 4},
            yticklabels={autres, {\texttt{A\_24\_P175612}}, {\texttt{A\_23\_P163278}}, {\texttt{A\_23\_P205428}}},
            nodes near coords,
            nodes near coords style={font=\tiny},
            tick label style={font=\tiny}
          ]
          \addplot[fill=geblue] coordinates {
              (0.08, 1)
              (0.16, 2)
              (0.27, 3)
              (0.81, 4)
            };
          \draw[red, dashed] (axis cs:0.6, 0.5) -- (axis cs:0.6, 4.5);
        \end{axis}
      \end{tikzpicture}
      \vskip 2pt
        {\footnotesize OvR Lasso : un gène robuste (\texttt{A\_23\_P205428}, $\hat{\Pi} \simeq 0{,}81$).}
    \end{column}

    \begin{column}{0.48\textwidth}
      \begin{itemize}
        \item un \textbf{noyau de gènes GE} est sélectionné de façon stable ;
        \item \textbf{SGCCA} : noyau \textit{plus} stable (3 gènes $> \pi_{\text{thr}}$, même gène de tête que le Lasso $\Rightarrow$ convergence inter-méthodes) ;
        \item l’\textbf{exclusion de CGH} par les méthodes data-driven est \textit{remarquablement stable} : un résultat robuste, non un artefact.
      \end{itemize}
    \end{column}
  \end{columns}
\end{frame}

% Slide 33
\begin{frame}{Le signal est-il réel ? Test de permutation}
  \vspace{14pt}
  \begin{columns}[T]
    \begin{column}{0.53\textwidth}
      \textbf{Test de permutation} (Ojala \& Garriga, 2010) :
      \begin{itemize}
        \item $H_0$ : aucune dépendance entre variables et étiquettes (BA attendue = 1/3) ;
        \item on \textit{permute} les labels, on ré-entraîne tout le pipeline, $B = 500$ fois ;
        \item \textbf{p-value} = proportion de permutations égalant ou dépassant le score observé.
      \end{itemize}

      \medskip
      \begin{block}{Résultat}
        Tous les modèles : \textbf{p-value} = 0,002. On rejette $H_0$ : le signal prédictif est statistiquement significatif et \textit{non dû au hasard}, malgré $p \gg n$.
      \end{block}
    \end{column}

    \begin{column}{0.43\textwidth}
      \begin{exampleblock}{Reproductibilité}
        Graines aléatoires fixées (plis, bootstrap, initialisation des solveurs) $\Rightarrow$ chiffres exactement reproductibles.
      \end{exampleblock}

      \begin{alertblock}{Robuste vs fragile}
        \textbf{Robuste} : signal porté par GE ; CGH non informatif ; \texttt{cort}/\texttt{dipg} bien séparées ; significativité globale.

        \medskip
        \textbf{Fragile} : classement fin des meilleures méthodes ; recall exact sur \texttt{midl} (se joue sur 1–2 patients) ; liste précise des gènes (noyau stable, périphérie variable).
      \end{alertblock}
    \end{column}
  \end{columns}
\end{frame}

% Slide 34: Section 6
\section{Travail bibliographique}

% Slide 35
\begin{frame}{Travail bibliographique : fondements et positionnement}
  Chaque choix méthodologique est ancré dans la littérature :
  \begin{description}
    \item[\textbf{Intégration multibloc supervisée}] RGCCA (Tenenhaus \& Tenenhaus, \textit{Psychometrika}, 2011) ; SGCCA / sélection de variables (Tenenhaus \textit{et al.}, \textit{Biostatistics}, 2014) ; \textit{block-relaxation} (De Leeuw, 1994 ; Tenenhaus \textit{et al.}, 2017 ; Girka \textit{et al.}, 2024 ; seuil : Gloaguen \textit{et al.}, 2017).
    \item[\textbf{Fusion coopérative}] Cooperative Learning multi-vues (Ding, Li, Narasimhan \& Tibshirani, \textit{PNAS}, 2022).
    \item[\textbf{Régularisation \& GLM}] Lasso / Elastic-Net et descente proximale ; logistique multinomiale (GLM, lien logit) ; pondération \texttt{class\_weight="balanced"}.
    \item[\textbf{Stabilité \& validation}] Bootstrap (Efron, 1979) ; sélection par stabilité (Meinshausen \& Bühlmann, 2010) ; test de permutation (Ojala \& Garriga, 2010).
  \end{description}

  \begin{columns}[T]
    \begin{column}{0.48\textwidth}
      \textbf{Filiation directe} : même jeu de données (gliomes IGR) et même cadre que le \textit{case study} RGCCA ; SGCCA + LDA \textbf{reproduit} la cible publiée (CV $\simeq 0{,}826 \pm 0{,}136$, test $\simeq 91{,}7\%$).
    \end{column}
    \begin{column}{0.48\textwidth}
      \textbf{Apports} : comparaison contrôlée des familles de modèles ; diagnostic de la non-convergence du Cooperative à $\rho > 0$ ; mécanisme OvR vs multinomial pour \texttt{midl} ; analyse de stabilité au-delà de la performance.
    \end{column}
  \end{columns}

  \bigskip
  \textit{Bibliographie complète dans le rapport écrit.}
\end{frame}

% Slide 36
\begin{frame}{Conclusion}
  \begin{block}{Ce que l’on retient}
    \begin{itemize}
      \item \textbf{Terrain difficile} : $p \gg n$ + multi-blocs + classe rare \texttt{midl} ;
      \item \textbf{SGCCA + LDA} : méthode retenue : la plus performante (test $\simeq 0{,}924$) \textit{et} la plus stable, avec recall \texttt{midl} 2/3 ;
      \item modèle \textbf{parcimonieux et interprétable} ($\simeq 68$ gènes, 11 régions CGH) ;
      \item signal validé par test de permutation ($p \simeq 0{,}002$).
    \end{itemize}
  \end{block}

  \medskip
  \textbf{Perspectives :}
  \begin{itemize}
    \item validation sur cohorte externe ; interprétation biologique des gènes stables ;
    \item plusieurs composantes SGCCA ; intégration d’un 3\ieme{} bloc (clinique).
  \end{itemize}

  \vfill
  \begin{center}
    {\large \textbf{Merci de votre attention \quad :)}}
  \end{center}
\end{frame}

\end{document}
